\section{An effective Hilbert--Samuel estimate on one component}
\label{sec:hs-component}
We derive the coefficient of \(\log(2+|h|)\) and the uniform additive
loss in \eqref{eq:cov-hs}. Fix a component of geometric degree
\(\delta\ge1\), height \(h\), and positive rank \(d\). Let its
unshifted and shifted arithmetic degrees be \(a_0,a_-\).
Write \(D\ge1\) for the flag dimension, \(T\ge3D+5\) for the length,
\(P\ge1\) for the size parameter and \(b\ge1\) for the error scale.
The hypotheses used below are
\begin{align*}
 &\log b\le P,\quad TP\le\sqrt b,\quad \rho,\kappa\ge0,\quad
 z=7(D^2/T+D\rho)\le\tfrac12,\quad zP\kappa\le b,\\
 &Q=\binom{T-1}{D}\delta,\qquad (1-z)d\le Q,\qquad
 a_-\ge a_0-dc,\quad c\ge0 .
\end{align*}
On this component the geometric comparison gives numbers
\(n,\Lambda,\beta,\Lambda_D,\Lambda_T\ge1\) and \(J\) such that
\begin{align*}
 &\log n,\log\beta,\log\delta\le P,\qquad
 \log\Lambda\le P(1+\kappa/T),\\
 &\log\Lambda_D,\log\Lambda_T\le P\kappa,\qquad
 J\le 4(DP(1+\kappa/T)+TP)+2DP\kappa/T .
\end{align*}
We spell out the geometric Hilbert--Samuel inequality used as input.
Set
\begin{align*}
 r&=\log\Lambda+26\log\big(nD\delta(2+|h|+\log\Lambda)\big),\\
 G&=(D+1)(\log n+1)/2+1,\\
 U&=Th-3(D+1)|h|-J-26D^2(r+\log T)
              -(T+D+1)\log\beta-(T+D^2)(G+2),\\
 V&=T(\log2+r+\tfrac12\log\beta)+\log\Lambda_D+\log\Lambda_T .
\end{align*}
The component assertion with constant \(26\) and initial multiplicity
one says \(0\le1-Q/d\le1\) and
\[
             a_0\ge Q U-(d-Q)V .
\]
It comes from an estimate for sections on the geometric component.
The following argument turns it into a degree bound with just one
height-dependent error and a constant independent of the component.

Put \(q=\kappa/T\) and \(l_0=27P(q+1)+65P\). Since \(T\le\sqrt b\)
we have \(\log D\le\log T\le P/2\). Also
\(\log(1+\log\Lambda)\le\log\Lambda\); hence
\[
 \log\Lambda+26\log\big(nD\delta(1+\log\Lambda)\big)\le l_0 .
\]
Indeed the three size logs contribute at most
\(26(P+P/2+P)=65P\), and the remaining terms at most
\(27\log\Lambda\). The inequality
\(2+|h|+\log\Lambda\le(2+|h|)(1+\log\Lambda)\) now yields
\(0\le r\le l_0+26\log(2+|h|)\).
From the rank share we have \(d\le2Q\) and \(d-Q\le zd\le2zQ\).
Since \(V\ge0\) the geometric comparison bounds \(a_0\) below by
\(Q(U-2zV)\). Substitution of the upper bound for \(r\) in its
negative terms gives
\begin{equation}
 a_-\ge Q\{Th-3(D+1)|h|-A\log(2+|h|)-H-2c\},\quad
       A=26(26D^2+2zT),                         \label{eq:hs-stage}
\end{equation}
where the complete constant \(H\) is
\begin{align*}
 H={}&J+26D^2(l_0+\log T)+(T+D+1)\log\beta\\
     &+(T+D^2)(G+2)
       +2z\{T(\log2+l_0+\tfrac12\log\beta)+\log\Lambda_D+\log\Lambda_T\}.
\end{align*}
In particular the rank defect multiplies the norm cost \(V\). Ignoring
it would lose the \(2zT\) term in \(A\).

Here are estimates for all five summands. The relations defining \(z\)
imply \(D^2\le zT/7\le T/14\), and therefore
\[
 D^2Pq\le b/7,\qquad D^2P\le b/14,\qquad
 DPq\le b/7,\qquad TP\le b,\qquad (TP)^2\le b .
\]
The first inequality follows by multiplying \(D^2\le zT/7\) by \(Pq\).
\begin{enumerate}
\item Expanding the bound for \(J\) gives \(6DPq+4DP+4TP\le9b\).
\item Since \(l_0+\log T\le27Pq+(185/2)P\), its contribution times
\(26D^2\) is at most \(26(27/7+185/28)b<300b\).
\item The logarithm of \(\beta\) contributes at most \(2TP\le2b\).
\item We have \(G+2\le5DP\) because \(P,D\ge1\) and \(\log n\le P\).
Also \(TD P\le (TP)^2\le b\) and \(D^2\le T\).
This gives at most \(10b\).
\item In the final term \(zTPq=zP\kappa\le b\). Both \(z\log\Lambda_D\)
and \(z\log\Lambda_T\) are bounded by \(b\). Moreover
\(zT\log2\le b/2\), \(zTP\le b/2\), and \(zT\log\beta/2\le b/4\).
Inserting \(l_0=27Pq+92P\) bounds the expression, including its
outer factor \(2\), by \(2(1/2+27+46+1/4+2)b<200b\).
\end{enumerate}
The five upper bounds give \(H\le4000b\).
Thus \eqref{eq:hs-stage} proves \eqref{eq:cov-hs}
with \(R=4000b+2c\).

There is also a cost for \(A\) in the subsequent component sum.
We have \(1\le A\le702T^2\), since \(D\le T\), \(2z\le1\) and \(T\ge1\).
Consequently \(\log A\le\log702+2\log T\le 8P\). Hence
\[
 A\log(\max(1,A))\le5616T^2P\le6000b .
\]
In the actual family the twist comparison supplies \(c=C_p Z_0^2\)
with \(C_p\ge0\) uniform and \(Z_0\) the base size. Its size conditions
give \(1\le Z_0,\ Z_0^{12}\le P,\ TP\le\sqrt b\), so \(Z_0^2\le b\).
As \(b\le E\), this proves
\(A\log(\max(1,A))+R\le(10000+2C_p)E\) with the same constant on the
common windows. The geometric component assertion and the twist
comparison, whose statements enter explicitly above, are the inputs
to this calculation.
\footnote{\href{https://github.com/Deicyde/riemann-hypothesis/blob/40e56bab/src/Formalism/LocalToGlobal/IndefiniteForms/IndefiniteTree/RealCubicMeans/LogMeanShares/LogMeanStatements/HSClaim.lean}{Component inequality};
\href{https://github.com/Deicyde/riemann-hypothesis/blob/40e56bab/src/Development/AbelianVarietyTheory/Proofs/DivisorProofs/DegBoundOfClaimNumbers.lean}{five errors and the twist};
\href{https://github.com/Deicyde/riemann-hypothesis/blob/40e56bab/src/Development/AbelianVarietyTheory/Proofs/DivisorProofs/HilbertSamuelALogALe.lean}{the logarithmic cost};
\href{https://github.com/Deicyde/riemann-hypothesis/blob/40e56bab/src/Development/AutomorphicArithmetic/ThetaCorrespondence/ThetaPairLifts/ThirdPlaceSeries/HilbertSamuelLowerBoundOfComponents.lean}{specialization at the common threshold}.}
